
%\needspace{4\baselineskip}
\color{uniblue}{\subsection{Контроль наличия напряжения на присоединении (КННп)}}
\color{black}

\begin{enumerate}

\item
В устройстве контроль наличия напряжения на присоединении представлен функцией «КННп». Функция контролирует наличие напряжения на отходящем присоединении. Функция используется для АПВ, а также для оперативного включения.

\item
Пусковые условия функции формируются без выдержки времени при превышении минимального или максимального из подведенных линейных напряжений выше уставки в зависимости от того сколько фаз будет подведено к устройству.

\item
Функция предусматривает два режима работы, определяемые уставкой SGF2:
\begin{itemize}
    \item измеренный;
    \item внешний.
\end{itemize}

Алгоритм работы функции, когда уставка SGF2 выставлена в положение «Измеренный»:
\begin{itemize}
    \item если к устройству подводятся три фазных или два линейных напряжения, то уставка SGF3 должна быть выставлена в положение «3 фазы». В этом случае пусковые условия будут формироваться, когда все три линейных напряжения будут превышать значение уставки с отсутствием срабатывания ИО обратной последовательности выше уставки;
    \item если к устройству подводится только одно линейное напряжение, то уставку SGF3 необходимо выставить в положение «2 фазы». В этом случае пусковые условия будут формироваться при превышении максимального из линейных напряжений выше уставки без учета срабатывания ИО обратной последовательности.
\end{itemize}

\item
Предусмотрена блокировка функции при срабатывании ИО обратной последовательности напряжения выше уставки.

\item
Если уставка SGF2 выставлена в положение «Внешний», то сигнал пуска формируется при наличии активного внешнего дискретного сигнала «Наличие Uп» без учета срабатывания измерительных органов.

\item
Ввод функции в работу на этапе параметрирования устройства осуществляется программным переключателем SGF1 «Ввод\_функции». Информация о введенной функции отображается выходным сигналом «КННп: Ввод».

\item
Функция может быть выведена из работы оперативно путем активации сигнала «ОВ КННп», что характеризуется активным состоянием сигнала «КННп: Оперативный вывод» на выходе функции. 

\item
Алгоритм работы функции выполнен в соответствии с рисунком \ref{knnp:img1}. В таблице \ref{knnp:tbl1} приведены параметры, необходимые для настройки функции.

\vspace{1mm}
\begin{figure}[H]
\centering
\includegraphics[width=1\textwidth,height=1\textheight,keepaspectratio]{\fbpath/123. КННп/_latex/img/knnp.pdf}
\captionof{figure}{Функциональная схема логики КННп}\label{knnp:img1}
\end{figure}


\end{enumerate}
